%% notes-tex/metabolic-module-report/metabolic-module-report.tex
%% [[experiments.026-metabolism-flux.enzyme-constrained-thermodynamic-flux-layer]]
%% https://github.com/Mjvolk3/torchcell/tree/main/notes-tex/metabolic-module-report/metabolic-module-report.tex
%%
%% A three-page digest of the metabolic module work: what the flux module is and what
%% feeds it, what the kinetic parameter tables it is meant to consume are worth, and
%% where the training results stand. Every figure and number is taken from the 026
%% kinetics-thermodynamics document, the 027 grid document and the flux-layer dev note,
%% and is produced by a committed script under experiments/026-metabolism-flux/ or
%% experiments/027-betaxanthin-metabolic/.
%%
%% Build:  make            -> metabolic-module-report.pdf
%%         make plots      -> re-convert the matplotlib panels from ASSET_IMAGES_DIR
%%         make check      -> formatting / provenance / style gate

\documentclass[10pt,a4paper]{article}
\usepackage[draft]{tcdoc}
\usepackage{float}
\usepackage{amsmath}
\usepackage{amssymb}

\emergencystretch=3.5em

%% Figures are placed with \tcfigh (tcdoc.sty), \tcfigfit placed [H] rather than floated:
%% on three pages a floated figure lands on its own page and leaves the text half empty.

\title{\vspace{-1.5em}The Metabolic Module in Three Pages\\
  {\large The flux layer, the parameters it is fed,\\and what has been measured}}
\author{Michael Volk}
\date{\today}

\begin{document}
\maketitle
\thispagestyle{empty}

\begin{abstract}
\noindent
The metabolic module is a differentiable layer that maps a perturbed gene representation to
a flux vector over yeast-GEM 9.0.2 under stoichiometry, an enzyme-capacity bound and a
second-law penalty. Its three parameter tables,
turnover numbers, Michaelis constants and formation energies, are meant to be tabulated
rather than learned. The tables have been built and audited: six published $k_{cat}$
predictors disagree on a single gene by a median factor of 19, yet 56 to 84 percent of
their values land inside the spread of enzymes sharing a function, against 53 percent for a
lookup that knows only the function. No trained arm consumes one yet, and the one training
grid run so far, 144 cells on Delta, produced constant predictors in every cell, so the
module's effect on betaxanthin prediction is unmeasured rather than negative. Full accounts
are in the 026 and 027 documents and the flux-layer note.
\end{abstract}

%% notes-tex/metabolic-module-report/sections/1-module.tex
%% [[torchcell.models.equivariant_cell_graph_transformer.mermaid.metabolic-module-simple]]
%% https://github.com/Mjvolk3/torchcell/tree/main/notes-tex/metabolic-module-report/sections/1-module.tex

\section{The flux module and what feeds it\secstatus{todo}}
\label{sec:module}

The cell graph transformer's instruments map representation to representation (Type~I) or
to a measured output (Type~II). The flux module is neither: it leaves gene index space for
reaction space $\mathbb{R}^{4131}$ under a fixed stoichiometry, and its output $\nu$ is a
latent physical state phenotypes are read from (Figure~\ref{fig:module}). A deletion sets a
gene's availability to zero as a hard fact; the GEM's gene-protein-reaction rules fold gene
gates to a reaction by a softmin over a complex's subunits and a sum over isozymes, and the
result scales an enzyme-capacity ceiling,
\begin{equation}
  |\nu_j| \le c_j(H_{\mathrm{pert}})\, k_{cat,j}\,\bar E,
  \qquad
  \mathcal{L} = \mathcal{L}_{\mathrm{data}}
  + \lambda_{\mathrm{th}}\,\frac{1}{|\mathcal{J}|}\sum_{j\in\mathcal{J}}
    \frac{[\,\nu_j\,\Delta_r G_j + \epsilon\,]_{+}}{|\nu_j|\,|\Delta_r G_j| + \epsilon}.
  \label{eq:module}
\end{equation}
The capacity bound is the only path by which a genotype reaches a flux magnitude, since
4{,}129 of the GEM's 4{,}131 bounds are $\pm1000$ or $0$; $\bar E$ is constant and no
enzyme abundance is predicted. The second-law term charges for flux running uphill and,
because $\Delta_r G$ sums to zero around any cycle, removes with no binary variable the 427
internal loops a growth constraint leaves free. It acts at training time only; the violation
fraction is reported at evaluation. Gibbs dissipation is a logged readout, not a
penalty: the Niebel 2019 limit is an empirical plateau, not a law.

\tcfigh[0.8]{figures/cgt_metabolic_module.pdf}{\textbf{The flux module's main components
  and what feeds them.} Orange, the path from the perturbed gene state to enzyme
  availability; purple, the fixed yeast-GEM 9.0.2 stoichiometry and bounds; yellow, the
  parameter tables, a \textbf{dashed} border marking one built and read by nothing; blue,
  the flux vector; red, the soft constraints. Three of the five tables are
  dashed: $k_{cat}$ resolves from the Open Enzyme
  Database at 3.7 percent of catalytic units and $\Delta_f G$ from the table shipped with
  the GEM. A predicted $k_{cat}$ tier and a merged energy route exist in the resolvers,
  under test. Source
  \texttt{notes/torchcell.models.equivariant\_cell\_graph\_transformer.mermaid.metabolic-module-simple.md}.}{fig:module}

%% notes-tex/metabolic-module-report/sections/2-parameters.tex
%% [[experiments.026-metabolism-flux.scripts.plot_per_gene_predictor_agreement]]
%% https://github.com/Mjvolk3/torchcell/tree/main/notes-tex/metabolic-module-report/sections/2-parameters.tex

\section{What the parameter tables are worth\secstatus{todo}}
\label{sec:parameters}

Six published $k_{cat}$ predictors (DLKcat, UniKP, EITLEM, TurNuP, DeepEnzyme, EZKinetics)
and four $K_M$ predictors were run from hash-pinned mirrors over 95.3 percent of yeast-GEM's
3{,}728 catalytic units, the ceiling set by substrate structure availability. Per gene, the
unit the capacity bound reads (Figure~\ref{fig:agreement}), the median gene moves 1.28
orders of magnitude between the highest and lowest predictor and 71 percent move by more
than one order, against a consensus range across genes of 0.87. On the 57 measured yeast
pairs, EZKinetics leads on rank at Spearman
0.79 and EITLEM on magnitude at 0.23 orders median error, an upper bound: every predictor
trained on the literature those measurements come from.

\tcfigh[0.86]{figures/per_gene_agreement_k_cat.pdf}{\textbf{Agreement between the six
  $k_{cat}$ predictors, and their accuracy against measured yeast values.} Panels a to d
  cover the 980 genes all six predictors cover; e and f the 57 measured pairs.
  \textbf{a}, Across-predictor spread per gene; dashed, the median of 1.28 orders of
  magnitude. \textbf{b}, The same spread against the dynamic range of $k_{cat}$, genes
  ranked by their median. \textbf{c}, Gene-level Spearman between predictors, 0.01 to 0.55.
  \textbf{d}, The betaxanthin supply genes; the shaded band marks the two heterologous
  cassette members with no yeast-GEM reaction. \textbf{e}, Predicted against measured; the
  band's half width of 0.75 is the median standard deviation of $\log_{10} k_{cat}$ within
  one EC4 class. \textbf{f}, Spearman and median absolute error, ranked. Figure 6 of the 026
  document, from \texttt{plot\_per\_gene\_predictor\_agreement.py}.}{fig:agreement}

\textbf{Most of that plausibility is coarse-class knowledge} (Figure~\ref{fig:plausibility}).
Against the EC4 band every predictor beats the EC-median lookup, 56 to 84 percent against
53; at EC1 the lookup already reaches 79 percent, and for $K_M$ no predictor beats the
lookup at all. On the nine unmeasured betaxanthin supply genes, four of six $k_{cat}$
predictors are class-plausible and no $K_M$ predictor is.

The thermodynamic table is recomputed with eQuilibrator component contribution, which
returns a covariance the shipped point values cannot. It prices 2{,}741 of 4{,}131
reactions, 986 recovered by fixing the proton convention; with the shipped table it reaches
3{,}716, 94.4 percent of flux-carrying reactions. \textbf{The shipped file is in kilocalories
and labeled kilojoules}, copied verbatim from yETFL by yeast-GEM's curation script:
regressing our table on yETFL's over 1{,}927 shared metabolites gives a slope of 4.187 at
$r = 0.987$. Correcting that and matching compartment conditions closes the transport arm of
the disagreement, median 17.4 to 5.6 kJ\,mol$^{-1}$, leaving the single-compartment cloud at
17.6, the genuine estimator difference. Of the 665 unpriced metabolites, 185 carry placeholder formulas and
are a method ceiling; \textbf{347 of the rest now carry a structure written from the name
and a price}, taking the recomputation past the shipped column to 86.3 percent, though the
union does not move and the values need re-transforming per compartment. All five
betaxanthin intermediates resolve by InChIKey, and Uni-pKa on its simple enumeration
template reproduces ChemAxon's energies to 0.91 kJ\,mol$^{-1}$, so the license binds
nowhere.

\tcfigh[0.86]{figures/within_class_plausibility_k_cat.pdf}{\textbf{How many of each
  predictor's $k_{cat}$ values fall inside the spread of enzymes doing the same job.}
  \textbf{a}, Fraction of the 57 measured pairs inside a shared band as the class is read
  to fewer EC digits; dashed, the EC-median lookup. \textbf{b}, The EC4 fraction against
  each pair's own class band, Wilson 95 percent intervals. \textbf{c}, The 23 betaxanthin
  supply units against one standard deviation of their EC4 class. \textbf{d}, Fraction of
  the 16 scorable units inside, by predictor: DLKcat 16 of 16, UniKP 15, EITLEM 6.
  \textbf{e}, Each unit's class band half width. Figure 7 of the 026 document, from
  \texttt{plot\_within\_class\_plausibility.py}.}{fig:plausibility}

%% notes-tex/metabolic-module-report/sections/3-status.tex
%% [[experiments.027-betaxanthin-metabolic]]
%% https://github.com/Mjvolk3/torchcell/tree/main/notes-tex/metabolic-module-report/sections/3-status.tex

\section{Where the training stands\secstatus{todo}}
\label{sec:status}

The 027 grid on Delta ran 144 cells, six arms over 24 seeds, on Cachera's betaxanthin
screen with the 019 nine-graph backbone. The anchored module scored $-0.020$ test Spearman
against $+0.027$ for no module. The prediction dumps show every cell is a constant
predictor: the largest prediction spread in any cell is 0.8 percent of the target
spread, the no-module arms sit at the target mean with a training loss of exactly 1.00 from
epoch 10, and the flux arms at wrong constants. The claim that the thermodynamic anchor
costs is therefore withdrawn and the module's effect is unmeasured; the same backbone
reached 0.434 under the 019 trainer, so the defect is in the 027 loop. Two defects precede
any relaunch: training uses yeast-GEM's shipped exchange bounds, which forbid all
amino-acid uptake, while the Cachera strains grew in a supplemented medium; and 170 of the
183 proton-translocating reactions carry a confidently wrong zero, the translocation term
being computed but off by default.

\textbf{Order of work.} A trainer that reproduces 0.434 on the no-module arm with an abort
on prediction collapse; the module arms on it, on the dataset's actual medium; then the
merged $\Delta_f G$ and predicted $k_{cat}$ tables, on GilaHyper jobs 1663, 1667, 1668 and
1669, carried to a selected epoch.


\end{document}
